%% ============================================================
%%  Chapter 1 — Introduction
%% ============================================================
\chapter{Introduction}
\label{ch:introduction}

\section{Motivation}

Cyber attacks have grown in sophistication and frequency, placing critical
infrastructure, financial systems, and personal data at persistent risk.
Traditional intrusion-detection systems (IDS) rely on static rule bases or
fixed-threshold anomaly detectors that struggle to adapt as attacker
strategies evolve.
Honeypots --- decoy systems that attract and study attackers --- offer a
complementary defence layer, but their effectiveness depends on when and how
a defender engages with each suspicious connection.

Reinforcement learning (RL) provides a natural framework for this decision
problem: an agent repeatedly chooses how to respond to incoming traffic and
receives a reward signal that encodes domain knowledge about correct
behaviour.
When combined with deep neural networks, RL agents can learn non-linear
policies from high-dimensional observations without hand-crafted feature
engineering.

\section{Problem Statement}

Two core challenges motivate this work.
First, realistic attacker simulation is difficult: a synthetic attacker must
produce plausible network traffic, follow a coherent escalation logic, and
support multiple distinct intent profiles.
Second, a learned honeypot policy must balance \emph{sensitivity}
(detecting real attacks) against \emph{specificity} (avoiding false alarms
that disrupt legitimate traffic), across a wide range of attack types and
kill-chain stages.

\section{Research Objectives}

This thesis addresses the following objectives:

\begin{enumerate}
  \item Design a Markov-chain attacker model grounded in the Lockheed Martin
        Cyber Kill Chain, with per-intent transition modifiers and
        UNSW-NB15-inspired feature distributions.
  \item Implement a DQN-based defender that learns an adaptive honeypot
        response policy from a 24-dimensional state vector.
  \item Develop a composite threat-level metric that integrates multiple
        attack signals and drives a domain-informed reward function.
  \item Evaluate the trained policy using detection rate, false-positive
        rate, episode reward, and DQN training loss across all four attacker
        intents.
\end{enumerate}

\section{Contributions}

The principal contributions of this work are:

\begin{itemize}
  \item A fully self-contained attacker--defender simulation environment
        built on the Gymnasium RL interface, enabling reproducible
        experiments without access to live network data.
  \item An intent-aware Markov attacker that generates qualitatively
        different attack campaigns from a shared base transition structure.
  \item A reward function that encodes honeypot domain knowledge ---
        including tarpitting bonuses, late kill-chain penalties, and
        attack-type-specific incentives --- as a learnable signal for DQN.
  \item Empirical evidence that a policy trained on \emph{Opportunistic}
        traffic generalises to harder intent profiles.
\end{itemize}

\section{Thesis Outline}

Chapter~\ref{ch:background} surveys the theoretical foundations: reinforcement
learning, Deep Q-Networks, Markov chains, the Cyber Kill Chain, honeypot
technology, and the UNSW-NB15 dataset.
Chapter~\ref{ch:methodology} describes the system design --- the attacker
module, defender module, Gymnasium environment, and evaluation infrastructure.
Chapter~\ref{ch:results} presents quantitative results from training and
cross-intent evaluation experiments.
Chapter~\ref{ch:discussion} interprets the findings, discusses limitations,
and identifies directions for future research.
Chapter~\ref{ch:conclusion} summarises the thesis.
